\documentclass[10pt]{article}  

%%%%%%%% PREÁMBULO %%%%%%%%%%%%
\title{Plantilla para reportes IMEC-UTB}
\usepackage[spanish]{babel} %Indica que escribiermos en español
\usepackage[utf8]{inputenc} %Indica qué codificación se está usando ISO-8859-1(latin1)  o utf8  
\usepackage{amsmath} % Comandos extras para matemáticas (cajas para ecuaciones,
% etc)
\usepackage{amssymb} % Simbolos matematicos (por lo tanto)
\usepackage{graphicx} % Incluir imágenes en LaTeX
\usepackage{color} % Para colorear texto
\usepackage{subfigure} % subfiguras
\usepackage{float} %Podemos usar el especificador [H] en las figuras para que se
% queden donde queramos
\usepackage{capt-of} % Permite usar etiquetas fuera de elementos flotantes
% (etiquetas de figuras)
\usepackage{sidecap} % Para poner el texto de las imágenes al lado
	\sidecaptionvpos{figure}{c} % Para que el texto se alinie al centro vertical
\usepackage{caption} % Para poder quitar numeracion de figuras
\usepackage{commath} % funcionalidades extras para diferenciales, integrales,
% etc (\od, \dif, etc)
\usepackage{cancel} % para cancelar expresiones (\cancelto{0}{x})
 
\usepackage{anysize} 					% Para personalizar el ancho de  los márgenes
\marginsize{2cm}{2cm}{2cm}{2cm} % Izquierda, derecha, arriba, abajo

\usepackage{appendix}
\renewcommand{\appendixname}{Apéndices}
\renewcommand{\appendixtocname}{Apéndices}
\renewcommand{\appendixpagename}{Apéndices} 

% Para que las referencias sean hipervínculos a las figuras o ecuaciones y
% aparezcan en color
\usepackage[colorlinks=true,plainpages=true,citecolor=blue,linkcolor=blue]{hyperref}
%\usepackage{hyperref} 
% Para agregar encabezado y pie de página
\usepackage{fancyhdr} 
\pagestyle{fancy}
\fancyhf{}
\fancyhead[L]{\footnotesize UTB} %encabezado izquierda
\fancyhead[R]{\footnotesize Fac. de Ingeniería}   % dereecha
\fancyfoot[R]{\footnotesize Reporte}  % Pie derecha
\fancyfoot[C]{\thepage}  % centro
\fancyfoot[L]{\footnotesize Ing. Mecanica y Mecatrónica}  %izquierda
\renewcommand{\footrulewidth}{0.4pt}


\usepackage{listings} % Para usar código fuente
\definecolor{dkgreen}{rgb}{0,0.6,0} % Definimos colores para usar en el código
\definecolor{gray}{rgb}{0.5,0.5,0.5} 
% configuración para el lenguaje que queramos utilizar
\lstset{language=Matlab,
   keywords={break,case,catch,continue,else,elseif,end,for,function,
      global,if,otherwise,persistent,return,switch,try,while},
   basicstyle=\ttfamily,
   keywordstyle=\color{blue},
   commentstyle=\color{red},
   stringstyle=\color{dkgreen},
   numbers=left,
   numberstyle=\tiny\color{gray},
   stepnumber=1,
   numbersep=10pt,
   backgroundcolor=\color{white},
   tabsize=4,
   showspaces=false,
   showstringspaces=false}

\newcommand{\sen}{\operatorname{\sen}}	% Definimos el comando \sen para el seno
%en español

\title{Plantilla para Reportes IMEC-UTB}
% Basada en la plantilla para reportes UPIITA de  Overleaf

%%%%%%%% TERMINA PREÁMBULO %%%%%%%%%%%%

\begin{document}

%%%%%%%%%%%%%%%%%%%%%%%%%%%%%%%%%% PORTADA %%%%%%%%%%%%%%%%%%%%%%%%%%%%%%%%%%%%%%%%%%%%
																					%%%
\begin{center}																		%%%
\newcommand{\HRule}{\rule{\linewidth}{0.5mm}}									%%%\left
 																					%%%
\begin{minipage}{0.48\textwidth} \begin{flushleft}
\includegraphics[scale = 0.25]{Imagenes/logo_2015_-_color_}
\end{flushleft}\end{minipage}
\begin{minipage}{0.48\textwidth} \begin{flushright}
\includegraphics[scale = 0.15]{Imagenes/Logo_UTB_sin_letras2}
\end{flushright}\end{minipage}

													 								%%%
\vspace*{1.0cm}								%%%
																					%%%	
\textsc{\huge Universidad Tecnológica \\ \vspace{5px} de Bolívar}\\[1.5cm]	

\textsc{\LARGE Programas de Ingeniería Mecánica \\ \vspace{5px} y Mecatrónica }\\[1.5cm]													%%%

\begin{minipage}{0.9\textwidth} 
\begin{center}																					%%%
\textsc{\LARGE Reporte EF 01 }
\end{center}
\end{minipage}\\[0.5cm]
%%%
    																				%%%
 			\vspace*{1cm}																		%%%
																					%%%
\HRule \\[0.4cm]																	%%%
{ \huge \bfseries Título específico de lo que ud. hizo}\\[0.4cm]	%%%
 																					%%%
\HRule \\[1.5cm]																	%%%
 																				%%%
																					%%%
\begin{minipage}{0.46\textwidth}													%%%
\begin{flushleft} \large															%%%

% Aqui a continuación pongan los nombres de los integrantes
\emph{Autores:}\\	
Juan Pérez Gados T000XXXYY\\
María Juana Pereira Galindo T000XXYYY\\
Autor 3\\ 
Autor 4
%%%
			%\vspace*{2cm}	
            													%%%
										 						%%%
\end{flushleft}																		%%%
\end{minipage}		
																%%%
\begin{minipage}{0.52\textwidth}		
\vspace{-0.6cm}											%%%
\begin{flushright} \large															%%%
\emph{Profesores:} \\																	%%%
Edgardo W. Arrieta Ortiz, MSc., IM.\\
Asesor 2	\\
Asesor 3														%%%
\end{flushright}																	%%%
\end{minipage}	
\vspace*{1cm}
%\begin{flushleft}
 	
%\end{flushleft}
%%%
 		\flushleft{\textbf{\Large Nombre del curso (IMEC XXZZ)}	}\\																		%%%
\vspace{2cm} 																				
\begin{center}	
Cartagena, Bol. D.T. y C. \\
{\large \today}																	%%%
 			\end{center}												  						
\end{center}							 											
																					
\newpage																		
%%%%%%%%%%%%%%%%%%%% TERMINA PORTADA %%%%%%%%%%%%%%%%%%%%%%%%%%%%%%%%

\tableofcontents 

\newpage

\section{Resumen.}



 Lo que se hizo, para qué sirve, cómo se hizo , qué resultados se obtuvieron, y porqué es importante haber logrado esto.  (de 120 palabras máximo).
 Ejemplo de cómo dar una cita bibliográfica \cite{IEEEreferencias:Ref1}.



\section{Introducción.} 

 Lo que Ud. escriba aquí debe responder a las preguntas hechas por cualquier lector en el mundo: 
 
\begin{itemize}
\item ¿De qué se trata este trabajo?
\item ¿Por qué estas personas hicieron este trabajo?
\item ¿Por qué vale la pena leerlo, y a quienes le puede interesar leerlo?
\item ¿Qué tiene de original o importante lo hecho aquí?
\item ¿Qué tan difícil fue lograr los resultados que se muestran aquí?
\item Puede incluir dibujos si los necesita para explicar y/o fotografías

\end{itemize}
Haec quae dicis admodum in omni tempore semper cómoi tu dicas dicere. Et ne forte dicant quod duplex est. Plurrimi maximus res est cogitans in quam solvitur quaestio satis perspicuum parcentes nulla divisio requiritur moralis accurate perpensis omnibus iudiciis et certa . Electus quisque valet mensura .
Haec quae dicis admodum in omni tempore semper cómoi tu dicas dicere. Et ne forte dicant quod duplex est. Plurrimi maximus res est cogitans in quam solvitur quaestio satis perspicuum parcentes nulla divisio requiritur moralis accurate perpensis omnibus iudiciis et certa . Electus quisque valet mensura .
Haec quae dicis admodum in omni tempore semper cómoi tu dicas dicere. Et ne forte dicant quod duplex est. Plurrimi maximus res est cogitans in quam solvitur quaestio satis perspicuum parcentes nulla divisio requiritur moralis accurate perpensis omnibus iudiciis et certa . Electus quisque valet mensura .
Haec quae dicis admodum in omni tempore semper cómoi tu dicas dicere. Et ne forte dicant quod duplex est. Plurrimi maximus res est cogitans in quam solvitur quaestio satis perspicuum parcentes nulla divisio requiritur moralis accurate perpensis omnibus iudiciis et certa . Electus quisque valet mensura .
Haec quae dicis admodum in omni tempore semper cómoi tu dicas dicere. Et ne forte dicant quod duplex est. Plurrimi maximus res est cogitans in quam solvitur quaestio satis perspicuum parcentes nulla divisio requiritur moralis accurate perpensis omnibus iudiciis et certa . Electus quisque valet mensura .

\cite{IEEEreferencias:Ref2}


\section{Información preliminar.}


No es el "Marco Teórico" que normalmente se llena con copiar ypegar de internet en trabajos de otros cursos. Por eso tiene otro nombre.

\begin{itemize}
\item  ¿ Qué otros trabajos o desarrollos similares se han hecho?
\item  ¿Cuál es la diferencia de este respecto a aquellos?
\item  ¿Qué modelos, teorías o cuerpos de conocimiento hay que tener presentes para leer el resto del trabajo presentado.
\item NO REPITA lo que otros ya escribieron y presentaron en documentos digitales (web/o textos). 

\item Descríba lo que se necesita para entender SU trabajo en sus PROPIAS palabras.
\item Necesariamente debe nombrar los documentos o autores a los que hace referencia (citelos) y sobretodo genere sus propios análisis al respecto.
\item Solo haga mención explícita de lo que Ud. VA A usar en su trabajo. 
\item Temas relacionados pero QUE NO se van a usar pueden mencionarse como citas.
\end{itemize}



Haec quae dicis admodum in omni tempore semper cómoi tu dicas dicere. Et ne forte dicant quod duplex est. Plurrimi maximus res est cogitans in quam solvitur quaestio satis perspicuum parcentes nulla divisio requiritur moralis accurate perpensis omnibus iudiciis et certa . Electus quisque valet mensura .
Haec quae dicis admodum in omni tempore semper cómoi tu dicas dicere. Et ne forte dicant quod duplex est. Plurrimi maximus res est cogitans in quam solvitur quaestio satis perspicuum parcentes nulla divisio requiritur moralis accurate perpensis omnibus iudiciis et certa . Electus quisque valet mensura .
Haec quae dicis admodum in omni tempore semper cómoi tu dicas dicere. Et ne forte dicant quod duplex est. Plurrimi maximus res est cogitans in quam solvitur quaestio satis perspicuum parcentes nulla divisio requiritur moralis accurate perpensis omnibus iudiciis et certa . Electus quisque valet mensura .


\section{Metodologías.}



\begin{itemize}
\item Diga cómo se hizo lo que ud. (o uds) hicieron.
\item  Sea claro y preciso.
\item  Si alguien quiere repetir su trabajo, debe poder usar esto como guía.
\item  Registre cada una de las iteraciones seguidas por sus procesos de diseño.
\item  Si hubo trabajo en equipo, aquí debe quedar especificada la participación de cada uno.

\item Aquí deben estar registradas (como mínimo) los siguientes :
\begin{enumerate}
\item  Diagramas de cuerpo libre,
\item  diagramas de cortante y momentos,
\item  diagramas de torsión,
\item  cálculos de velocidades,
\item  análisis de aceleraciones y
\item  análisis de fuerzas,
\item  análisis de fatiga si es el caso,
\item  calculo de los factores de Marin y
\item  uso de variables estocásticas,
\item análisis de pandeo,
\item  diseño de ejes,
\item  diseño de sujetadores y uniones si los hay, etc
\end{enumerate}



\end{itemize}
Haec quae dicis admodum in omni tempore semper cómoi tu dicas dicere. Et ne forte dicant quod duplex est. Plurrimi maximus res est cogitans in quam solvitur quaestio satis perspicuum parcentes nulla divisio requiritur moralis accurate perpensis omnibus iudiciis et certa . Electus quisque valet mensura .
Haec quae dicis admodum in omni tempore semper cómoi tu dicas dicere. Et ne forte dicant quod duplex est. Plurrimi maximus res est cogitans in quam solvitur quaestio satis perspicuum parcentes nulla divisio requiritur moralis accurate perpensis omnibus iudiciis et certa . Electus quisque valet mensura .
Haec quae dicis admodum in omni tempore semper cómoi tu dicas dicere. Et ne forte dicant quod duplex est. Plurrimi maximus res est cogitans in quam solvitur quaestio satis perspicuum parcentes nulla divisio requiritur moralis accurate perpensis omnibus iudiciis et certa . Electus quisque valet mensura .




\begin{figure}[H]
	\begin{center}
 		\includegraphics[width = 0.5\textwidth]{Imagenes/silla3ruedas}
 		\captionof{figure}{\label{fig:silla3ruedas}Silla de ruedas para subir escaleras} 
	\end{center} 
\end{figure}



\subsection{Resultados.}


Describa los resultados obtenidos. 
Valores de Fuerzas, Cargas, 
Valores de esfuerzos hallados, 
si es el caso las dimensiones calculadas, espesores, longitudes, y 
planos de diseño resultantes.


Haec quae dicis admodum in omni tempore semper cómoi tu dicas dicere. Et ne forte dicant quod duplex est. Plurrimi maximus res est cogitans in quam solvitur quaestio satis perspicuum parcentes nulla divisio requiritur moralis accurate perpensis omnibus iudiciis et certa . Electus quisque valet mensura .
Haec quae dicis admodum in omni tempore semper cómoi tu dicas dicere. Et ne forte dicant quod duplex est. Plurrimi maximus res est cogitans in quam solvitur quaestio satis perspicuum parcentes nulla divisio requiritur moralis accurate perpensis omnibus iudiciis et certa . Electus quisque valet mensura .
Haec quae dicis admodum in omni tempore semper cómoi tu dicas dicere. Et ne forte dicant quod duplex est. Plurrimi maximus res est cogitans in quam solvitur quaestio satis perspicuum parcentes nulla divisio requiritur moralis accurate perpensis omnibus iudiciis et certa . Electus quisque valet mensura .

Haec quae dicis admodum in omni tempore semper cómoi tu dicas dicere. Et ne forte dicant quod duplex est. Plurrimi maximus res est cogitans in quam solvitur quaestio satis perspicuum parcentes nulla divisio requiritur moralis accurate perpensis omnibus iudiciis et certa . Electus quisque valet mensura .
Haec quae dicis admodum in omni tempore semper cómoi tu dicas dicere. Et ne forte dicant quod duplex est. Plurrimi maximus res est cogitans in quam solvitur quaestio satis perspicuum parcentes nulla divisio requiritur moralis accurate perpensis omnibus iudiciis et certa . Electus quisque valet mensura .




\section{Conclusiones.}

Explique por qué su diseño sirve (o nó) y qué hace falta por hacer.
Responda las expectativas generadas en la Introducción.
Analice sus resultados críticamente y diga si son creibles y responda
qué tan seguras son sus propuestas. 



%%%%%%% Bibliografía %%%%%%%%
\bibliographystyle{bst/IEEEtran.bst} 
\addcontentsline{toc}{section}{Referencias}  
\bibliography{bib/IEEEabrv,bib/IEEEreferencias.bib} 
%%%%%%% Bibliografía %%%%%%%%    

\appendix  
\clearpage % o \cleardoublepage
\addappheadtotoc 
\appendixpage 

\section{Memoria de cálculos.}

Digitalice sus manuscritos y los procedimientos de cálculo utilizados. Deben ser legibles y claros.

< Documentos escaneados> con algunas explicaciones.

Si prefiere escribir todos sus calculos usando fórmulas de LATEX hágalo a través de la metodología.

\section{Algoritmos.}

 Anexe copia en texto aquí y  de los algoritmos y demás archivos de soportes utilizados (documentos de Scilab, Maxima, NX, etc...).
 
 Los mismos archivos deben ser suplidos en forma digital en un archivo comprimido acompañando la entrea de este PDF.
 
 Aqui puede colocar capturas de pantalla de sus modelos en NX. 
 
 
 \section{Planos}
Los planos del deiseño presentado deben estar aquí. Generados por NX/Solid-edge.

Haec quae dicis admodum in omni tempore semper cómoi tu dicas dicere. Et ne forte dicant quod duplex est. Plurrimi maximus res est cogitans in quam solvitur quaestio satis perspicuum parcentes nulla divisio requiritur moralis accurate perpensis omnibus iudiciis et certa . Electus quisque valet mensura .
Haec quae dicis admodum in omni tempore semper cómoi tu dicas dicere. Et ne forte dicant quod duplex est. Plurrimi maximus res est cogitans in quam solvitur quaestio satis perspicuum parcentes nulla divisio requiritur moralis accurate perpensis omnibus iudiciis et certa . Electus quisque valet mensura .
Haec quae dicis admodum in omni tempore semper cómoi tu dicas dicere. Et ne forte dicant quod duplex est. Plurrimi maximus res est cogitans in quam solvitur quaestio satis perspicuum parcentes nulla divisio requiritur moralis accurate perpensis omnibus iudiciis et certa . Electus quisque valet mensura .




\end{document}